\specialchapter{\texorpdfstring{图版 VI. \(E_{7}\) 型根系}{图版 VI. E\_7 型根系}}

\begin{enumerate}
\def\labelenumi{(\Roman{enumi})}
\tightlist
\item
  V 是 \(E = R^{8}\) 中正交于 \(\varepsilon_{7} + \varepsilon_{8}\) 的超平面。
\end{enumerate}

根：\(\pm \varepsilon_{i} \pm \varepsilon_{j} (1 \leqslant i < j \leqslant 6)\)，\(\pm (\varepsilon_{7} - \varepsilon_{8})\)，

\[
\pm \frac {1}{2} (\varepsilon_ {7} - \varepsilon_ {8} + \sum_ {i = 1} ^ {6} (- 1) ^ {\nu (i)} \varepsilon_ {i}) \quad \text { with } \sum_ {i = 1} ^ {8} \nu (i) \text { odd. }
\]

根的数目：n = 126。

\begin{enumerate}
\def\labelenumi{(\Roman{enumi})}
\setcounter{enumi}{1}
\tightlist
\item
  基：
\end{enumerate}

\[
\begin{array}{r l} \alpha_ {1} & = \frac {1}{2} (\varepsilon_ {1} + \varepsilon_ {8}) - \frac {1}{2} (\varepsilon_ {2} + \varepsilon_ {3} + \varepsilon_ {4} + \varepsilon_ {5} + \varepsilon_ {6} + \varepsilon_ {7}), \\ \alpha_ {2} & = \varepsilon_ {1} + \varepsilon_ {2}, \alpha_ {3} = \varepsilon_ {2} - \varepsilon_ {1}, \alpha_ {4} = \varepsilon_ {3} - \varepsilon_ {2}, \alpha_ {5} = \varepsilon_ {4} - \varepsilon_ {3}, \\ \alpha_ {6} & = \varepsilon_ {5} - \varepsilon_ {4}, \alpha_ {7} = \varepsilon_ {6} - \varepsilon_ {5}. \end{array}
\]

正根：

\[
\begin{array}{c} \pm \varepsilon_ {i} + \varepsilon_ {j} (1 \leqslant i <   j \leqslant 6), - \varepsilon_ {7} + \varepsilon_ {8}, \\ \frac {1}{2} (- \varepsilon_ {7} + \varepsilon_ {8} + \sum_ {i = 1} ^ {6} (- 1) ^ {\nu (i)} \varepsilon_ {i}) \quad \text { with } \sum_ {i = 1} ^ {6} \nu (i) \text { odd }. \end{array}
\]

含有 \(\alpha_{7}\) 且至少有一个系数 \(\geqslant 2\) 的正根（我们把根 \(a\alpha_{1} + b\alpha_{2} + c\alpha_{3} + d\alpha_{4} + e\alpha_{5} + f\alpha_{6} + g\alpha_{7}\) 记为 \(\frac{a c d e f g}{b})^{14}\)：

\[
\begin{array}{c c c c c} 0   1   2   1   1   1 & 1   1   2   1   1   1 & 0   1   2   2   1   1 & 1   2   2   1   1   1 & 1   1   2   2   1   1 \\ 1 & 1 & 1 & 1 & 1 \\ 0   1   2   2   2   1 & 1   2   2   2   1   1 & 1   1   2   2   2   1 & 1   2   2   2   2   1 & 1   2   3   2   1   1 \\ 1 & 1 & 1 & 1 & 1 \\ 1   2   3   2   2   1 & 1   2   3   2   1   1 & 1   2   3   3   2   1 & 1   2   3   2   2   1 & 1   2   3   3   2   1 \\ 1 & 2 & 1 & 2 & 2 \end{array}
\]

\[
\begin{array}{c c c} 1 2 4 3 2 1 & 1 3 4 3 2 1 & 2 3 4 3 2 1 \\ 2 & 2 & 2 \end{array}
\]

\begin{enumerate}
\def\labelenumi{(\Roman{enumi})}
\setcounter{enumi}{2}
\item
  Coxeter 数：\(h = 18\) 。
\item
  最高根：
\end{enumerate}

\[
\tilde {\alpha} = \varepsilon_ {8} - \varepsilon_ {7} = 2 \alpha_ {1} + 2 \alpha_ {2} + 3 \alpha_ {3} + 4 \alpha_ {4} + 3 \alpha_ {5} + 2 \alpha_ {6} + \alpha_ {7} = \omega_ {1}.
\]

扩展 Dynkin 图：

\[
\alpha_ {1} \quad \alpha_ {3} \quad \alpha_ {4} \quad \alpha_ {5} \quad \alpha_ {6} \quad \alpha_ {7}
\]

\begin{enumerate}
\def\labelenumi{(\Alph{enumi})}
\setcounter{enumi}{21}
\tightlist
\item
  \(R^{\prime} = R.\)
\end{enumerate}

\[
  \Phi_ {\mathrm{R}} (x, y) = (x | y) / 36, \quad \gamma (\mathrm{R}) = 2 ^ {2}. 3 ^ {4}.
\]

\begin{enumerate}
\def\labelenumi{(\Roman{enumi})}
\setcounter{enumi}{5}
\tightlist
\item
  基本权：
\end{enumerate}

\[
\begin{array}{l} \omega_ {1} = \varepsilon_ {8} - \varepsilon_ {7} = 2 \alpha_ {1} + 2 \alpha_ {2} + 3 \alpha_ {3} + 4 \alpha_ {4} + 3 \alpha_ {5} + 2 \alpha_ {6} + \alpha_ {7} \\ \omega_ {2} = \frac {1}{2} (\varepsilon_ {1} + \varepsilon_ {2} + \varepsilon_ {3} + \varepsilon_ {4} + \varepsilon_ {5} + \varepsilon_ {6} - 2 \varepsilon_ {7} + 2 \varepsilon_ {8}) \\ \qquad = \frac {1}{2} (4 \alpha_ {1} + 7 \alpha_ {2} + 8 \alpha_ {3} + 12 \alpha_ {4} + 9 \alpha_ {5} + 8 \alpha_ {6} + 3 \alpha_ {7}) \\ \omega_ {3} = \frac {1}{2} (- \varepsilon_ {1} + \varepsilon_ {2} + \varepsilon_ {3} + \varepsilon_ {4} + \varepsilon_ {5} + \varepsilon_ {6} - 3 \varepsilon_ {7} + 3 \varepsilon_ {8}) \\ \qquad = 3 \alpha_ {1} + 4 \alpha_ {2} + 6 \alpha_ {3} + 8 \alpha_ {4} + 6 \alpha_ {5} + 4 \alpha_ {6} + 2 \alpha_ {7} \\ \omega_ {4} = \varepsilon_ {3} + \varepsilon_ {4} + \varepsilon_ {5} + \varepsilon_ {6} + 2 (\varepsilon_ {8} - \varepsilon_ {7}) \\ \qquad = 4 \alpha_ {1} + 6 \alpha_ {2} + 8 \alpha_ {3} + 12 \alpha_ {4} + 9 \alpha_ {5} + 6 \alpha_ {6} + 3 \alpha_ {7} \\ \omega_ {5} = \varepsilon_ {4} + \varepsilon_ {5} + \varepsilon_ {6} + \frac {3}{2} (\varepsilon_ {8} - \varepsilon_ {7}) \\ \qquad = \frac {1}{2} (6 \alpha_ {1} + 9 \alpha_ {2} + 12 \alpha_ {3} + 18 \alpha_ {4} + 15 \alpha_ {5} + 10 \alpha_ {6} + 5 \alpha_ {7}) \\ \omega_ {6} = \varepsilon_ {5} + \varepsilon_ {6} - \varepsilon_ {7} + \varepsilon_ {8} \\ \qquad = 2 \alpha_ {1} + 3 \alpha_ {2} + 4 \alpha_ {3} + 6 \alpha_ {4} + 5 \alpha_ {5} + 4 \alpha_ {6} + 2 \alpha_ {7} \\ \omega_ {7} = \varepsilon_ {6} + \frac {1}{2} (\varepsilon_ {8} - \varepsilon_ {7}) \\ \qquad = \frac {1}{2} (2 \alpha_ {1} + 3 \alpha_ {2} + 4 \alpha_ {3} + 6 \alpha_ {4} + 5 \alpha_ {5} + 4 \alpha_ {6} + 3 \alpha_ {7}). \end{array}
\]

\begin{enumerate}
\def\labelenumi{(\Roman{enumi})}
\setcounter{enumi}{6}
\tightlist
\item
  正根之和：
\end{enumerate}

\[
\begin{array}{l} 2 \rho = 2 \varepsilon_ {2} + 4 \varepsilon_ {3} + 6 \varepsilon_ {4} + 8 \varepsilon_ {5} + 10 \varepsilon_ {6} - 17 \varepsilon_ {7} + 17 \varepsilon_ {8} \\ = 34 \alpha_ {1} + 49 \alpha_ {2} + 66 \alpha_ {3} + 96 \alpha_ {4} + 75 \alpha_ {5} + 52 \alpha_ {6} + 27 \alpha_ {7}. \end{array}
\]


\begin{enumerate}
\def\labelenumi{(\Roman{enumi})}
\setcounter{enumi}{7}
\item
  P(R)/Q(R) 同构于 Z/2Z。连接指标：2。
\item
  指数：1, 5, 7, 9, 11, 13, 17。
\item
  W(R) 的阶：\(2^{10}.3^{4}.5.7\) 。
\item
  A(R) = W(R), \(w_{0} = -1\) .
\item
  \(P(R^{\prime})/Q(R^{\prime})\) 只有一个非恒等元：它确定 Dynkin 图的唯一非平凡自同构。
\end{enumerate}

\[
\left( \begin{array}{c c c c c c c} 2 & 0 & - 1 & 0 & 0 & 0 & 0 \\ 0 & 2 & 0 & - 1 & 0 & 0 & 0 \\ - 1 & 0 & 2 & - 1 & 0 & 0 & 0 \\ 0 & - 1 & - 1 & 2 & - 1 & 0 & 0 \\ 0 & 0 & 0 & - 1 & 2 & - 1 & 0 \\ 0 & 0 & 0 & 0 & - 1 & 2 & - 1 \\ 0 & 0 & 0 & 0 & 0 & - 1 & 2 \end{array} \right)
\]

\begin{enumerate}
\def\labelenumi{(\Roman{enumi})}
\setcounter{enumi}{12}
\tightlist
\item
  Cartan 矩阵：
\end{enumerate}

